\documentclass[12pt]{article}
\usepackage{preproposal}

\begin{document}
\thesisheader{Large Language Models as a Second Pass\texorpdfstring{\\}{ }%
  for Quantum Circuit Optimization}{Honours Thesis Preproposal}

\section{Introduction}

Test whether a large language model (LLM) can reduce the number of controlled-NOT gates (CX gates,
which flip a target qubit when a control qubit is 1) in a circuit already optimized by Qiskit
level 3 (the highest built-in optimization preset in IBM's compiler). A shorter circuit will count
only if it performs the same operation on every input. The primary comparison is the best verified
result from an LLM against the best verified result from two additional conventional optimizers,
starting from the same compiled circuit and given the same amount of extra time~\cite{ibm-stages}.

\researchquestion{When can LLMs improve an already optimized quantum circuit?}

\section{Motivation and Related Work}

Gu et al.\ show that a language model can design quantum circuits from a task specification.
StabilizerBench shows that several models can shorten existing circuits that prepare a required
state from zero. Our circuits act as reusable subroutines (operations called inside a larger
computation), so a rewrite must work on any input, including a superposition (a combination of
quantum states). Both precedents motivate giving the models a circuit task with an independent
correctness check~\cite{gu,stabilizerbench}.

Two-qubit gates are a frequent source of device errors, which makes CX count a useful
compiler-level measure. Conventional optimizers can already reduce gate counts further after
Qiskit level 3. We give those methods and the models the same compiled gate list and the same
additional time to ask which route finds the shorter verified circuit~\cite{ibm-opt,quspect}.

The working hypothesis is that additional verified savings will concentrate in specific circuit
constructions. In particular, a model could spot a better way to organize the carry operations
(the steps that move overflow from one digit to the next) of a ripple-carry adder after those
operations have been expanded into low-level gates. Other circuit groups give the hypothesis a
meaningful test. We will also check whether the models reach better final results when they see
the same circuits before level 3.

The conclusions will concern the tested circuit sizes, gate basis, and model configurations.
Gate-count gains are compiler-level results. Hardware performance needs its own device-specific
measurements.

\Needspace{17\baselineskip}
\section{Project Scope}

Use 48 fixed unitary circuits (circuits that perform reversible quantum operations) on 3--8
qubits, with 16 circuits in each group. Keep the original qubit count and wire order throughout.
The workload includes several constructions so that one adder or one layout cannot stand in for
its entire group.

\par\smallskip
\begingroup\small\noindent
\begin{tabularx}{\textwidth}{@{}
  >{\raggedright\arraybackslash}p{0.175\textwidth}
  >{\raggedright\arraybackslash}X
  >{\raggedright\arraybackslash}p{0.105\textwidth}
  @{}}
  \toprule
  \textbf{Group} & \textbf{Construction} & \textbf{Cases} \\
  \midrule
  Arithmetic &
    Qiskit's Cuccaro ripple-carry adder (\texttt{adder\_ripple\_c04}), with register widths 1--3
    and all three modes. The VBE ripple-carry adder (\texttt{adder\_ripple\_v95}), with width 2
    and all three modes. The Draper adder (\texttt{adder\_qft\_d00}, which adds by shifting
    phases in a Fourier representation of the numbers), with widths 2 and 3 in fixed and half
    modes. &
    9 + 3 + 4 \\
  \addlinespace[2pt]
  Simulation &
    One first-order Trotter step (apply the individual interactions in sequence) for each of two
    nearest-neighbor Hamiltonians (operators specifying the interactions of a quantum system),
    on chains of 3--6 qubits, at two evolution times. &
    4 \(\times\) 2 \(\times\) 2 \\
  \addlinespace[2pt]
  Variational circuits &
    Qiskit's \texttt{real\_amplitudes} ansatz (a circuit template with adjustable rotation
    angles), on 3--6 qubits, with linear or circular entanglement (CX gates between neighboring
    qubits on a chain or ring) and one or two repetitions. Bind the angles (set their numerical
    values) before compilation. &
    4 \(\times\) 2 \(\times\) 2 \\
  \bottomrule
\end{tabularx}
\par\endgroup\smallskip

In the arithmetic group, use the constructors from \texttt{qiskit.synthesis.arithmetic.adders}.
Width is the number of qubits in each input register. Fixed mode discards overflow, half mode
writes an additional carry-out, and full mode has both a carry-in and a carry-out. The circuit to
be preserved includes all declared wires, including any ancillas (temporary working qubits). The
benchmark checks every input on those wires, including any scratch inputs that a typical caller
initializes to zero~\cite{cuccaro,vbe,draper}.

For simulation, use the interactions below with unit coefficients. $X$, $Y$, and $Z$ are Pauli
gates (three basic one-qubit operations), and a product such as $Z_i Z_{i+1}$ acts on both
indicated wires. Implement each term $P$ once as $e^{-itP}$ (rotation generated by that
interaction) using Qiskit's \texttt{PauliEvolutionGate} with a single Pauli term. The circuit
produced by the given term order is the object to be preserved.

\begin{enumerate}
  \item Apply each $Z_i Z_{i+1}$ interaction, then each $X_i$ field.
  \item Apply each $X_i X_{i+1}$ interaction, then each $Y_i Y_{i+1}$, then each $Z_i Z_{i+1}$,
  then each $Z_i$ field.
\end{enumerate}

Within each set of terms, traverse the wires or bonds from left to right. Enumerate sizes in
ascending order, then Hamiltonians in the order above, then times $t=\pi/16,\pi/8$. For the
variational circuits, enumerate sizes in ascending order, then linear before circular, then one
before two repetitions. Use Python's \texttt{random.Random(20261124)} to draw angles in Qiskit's
parameter order from the ordered list $[-3\pi/8,-\pi/4,-\pi/8,\pi/8,\pi/4,3\pi/8]$.

\Needspace{10\baselineskip}
\section{Experimental Methodology}
\subsection{Development and Controls}

Develop and debug on six other instances. Use width-1 VBE fixed and full adders, the first
Hamiltonian on 3 and 4 qubits at $t=\pi/12$, and 3- and 4-qubit real-amplitudes circuits with one
\texttt{reverse\_linear} repetition and all angles $\pi/12$. Use separate two-qubit sanity checks
(cases with known answers), including identity (no change), CX, iSWAP (a swap with phase factors),
and SWAP, with known minima of 0, 1, 2, and 3 CX. Freeze the evaluation circuits, scoring code,
prompts, model configurations, run order, and baseline settings in a versioned manifest (a file
recording the exact experiment inputs) before collecting evaluation outputs~\cite{shende}.

\subsection{Starting Points and Methods}

Compile each original circuit with Qiskit 2.6.2 to a common basis (the allowed gate set) of
\texttt{u} (an arbitrary one-qubit gate) and \texttt{cx}. Use level 0 to lower composite
operations (expand named operations into the common gate set). This produces \textbf{Before level
3}. Apply level 3 to that circuit to produce \textbf{After level 3}. Both calls use
\texttt{seed\_transpiler=17}, \texttt{approximation\_degree=1.0},
\texttt{qubits\_initially\_zero=False}, and no device or coupling map. Any pair of logical qubits
(numbered wires in the circuit) can interact. Restore declared wire order with explicit gates if a
tool reports a final permutation (a reordering of the wires), and charge those gates to its
result. Verify both starts before using them~\cite{ibm-api}.

\par\smallskip
\begingroup\small\noindent
\begin{tabularx}{\textwidth}{@{}
  >{\raggedright\arraybackslash}p{0.228\textwidth}
  >{\raggedright\arraybackslash}X
  @{}}
  \toprule
  \textbf{Method} & \textbf{Procedure from a given starting circuit} \\
  \midrule
  Starting circuit &
    Keep the verified input. The level-3 version is the ordinary compiler result. \\
  \addlinespace[2pt]
  Extra conventional optimization &
    Run tket 2.18.5 with pytket-qiskit 0.79.0. Use
    \texttt{FullPeepholeOptimise(allow\_swaps=False, target\_2qb\_gate=OpType.CX)} followed by
    rebasing to TK1/CX (expressing the result with arbitrary one-qubit gates and CX), with
    implicit swaps disabled. Give tket at most 10 seconds. Use the remainder of the 180-second
    total for QUESO, independently starting from the same input. \\
  \addlinespace[2pt]
  GPT-5.6-Sol &
    Use the signed-in Codex CLI in text-only mode with medium reasoning effort. Start each call
    in an empty temporary directory with tools and project context disabled. Accept a candidate
    only from calls with no tool events. \\
  \addlinespace[2pt]
  Qwen3-Coder 30B &
    Use the Ollama \texttt{qwen3-coder:30b} build with Q4\_K\_M quantization (compressed model
    weights) and temperature 0.3 (a sampling setting). Allow 12,288 tokens (chunks of text) in
    context and 3,584 per reply. Use seeds 100 and 101 for the two submissions, and 102 and 103
    on repeat runs. \\
  \bottomrule
\end{tabularx}
\par\endgroup\smallskip

Run QUESO at GUOQ revision \texttt{8c4c3a5a6dfc9f7fc375ec16c2180139f0a8cb1a},
with Nam gates (\texttt{cx,h,rz,x}) and \texttt{-g NAM -opt TWO\_Q -search BEAM -temp 0 -q 8000
-resynth NONE}. Translate the shared input into Nam gates with Qiskit level 0 and verify that the
adapter preserves the operation and CX count. Normalize each completed candidate back to the
shared basis for scoring. Log conventional results separately as well as the best result obtained
by the two-tool sequence. Archive both model identifiers, provider configurations, hardware,
dependency versions, and Ollama's model digest (the identifier for the exact model files) in the
manifest~\cite{tket-adapter,guoq-code}.

\Needspace{5\baselineskip}
The two models represent a hosted reasoning system and a locally runnable coding model. Confirm
access on the development cases and record any required change to the lineup before freezing the
evaluation. Interpret differences as results for those configured systems, which also differ in
hosting and decoding (how model replies are generated).

\subsection{Model Interaction and Time Budget}

Give each model the qubit count, the input CX count, and the gate list in JSON (structured text).
Use the same rendering rule and task instructions for both starting points. Render standard
one-qubit gates as named aliases when the matrix matches up to phase within $10^{-12}$, with a
fixed lookup order of H, X, Y, Z, S, S\textdagger{}, T, and T\textdagger{}. Render other one-qubit
gates as \texttt{u}. Express angles as decimal multiples of $\pi$ rounded to 15 decimal places
when serializing. For example, \texttt{["rz", 0, 0.25]} denotes a rotation around Z by $\pi/4$ on
wire 0. Verify the serialized inputs against the original circuit. The prompt supplies the same
fields for all cases, with the model inferring any structure from the gates.

Each trial has at most two submissions and a hard 180-second wall-clock limit (actual elapsed
time). The model submits an ordered JSON array with CX and one-qubit gates. The allowed one-qubit
gates are \texttt{u,h,x,y,z,s,sdg,t,tdg,rx,ry,rz,p}. Accept a plain array or a single fenced JSON
block and enforce valid gate names, wire indices, and finite angles. After the first submission,
give the model the verifier status and the candidate's CX count if it parsed, plus a format error
if applicable. The second submission receives only that trial's conversation. Each model and
starting point gets an independent trial.

For each extra method, start the clock before input conversion or prompt serialization and include
inference (generating the model's reply) or search, provider wait, parsing, and verification. At
180 seconds, use only verified candidates already available. QUESO can submit saved intermediate
results as its search runs. Keep the starting circuit available throughout. Run methods serially
with one active trial at a time. Shuffle the main model trials with Python's
\texttt{random.Random(20261125)} and save the order in the manifest. Log time to the best accepted
CX count, overall elapsed time, and failure reason. Equal wall-clock limits make the comparison
about waiting time for the configured systems. The systems can use different hardware and amounts
of computation.

Run all 48 circuits at both starting points with both models, for 192 main trials and at most 384
submissions. Repeat both starts with both models on 12 prespecified circuits. Use Cuccaro width-2
fixed and width-3 full, VBE width-2 half, and Draper width-3 half, plus positions 1, 6, 11, and 16
in each of the ordered simulation and variational groups. This adds 48 repeat trials and at most
96 submissions. Keep repeats separate when reporting the main 48-circuit results. The planned cap
for evaluation inference is 480 submissions and 12 hours of serial model-trial time, plus up to
4.8 hours for conventional searches across the two starts.

\Needspace{10\baselineskip}
\subsection{Verification and Scoring}

Verify each input and candidate against the original circuit's unitary (its full matrix of action
on quantum states) in the original wire order. Let $U$ be the original matrix and $V$ the
candidate. Compute $s=\operatorname{Tr}(V^\dagger U)$ (a complex measure of their overlap) and use
the global phase factor $\lambda=s/|s|$ when $s$ is nonzero, or $\lambda=1$ otherwise. Require the
maximum absolute entry of $U-\lambda V$ to be at most $10^{-8}$ in double precision. This compares
the full operation up to a single common phase (a factor shared by every component of the quantum
state) and numerical tolerance. The convention covers standalone circuits on the given wires.
Coherently controlling a subroutine (applying it on one branch of a quantum control) would require
tracking that phase.

\Needspace{8\baselineskip}
Reject candidates that change the width or contain measurements, resets, classical control, or
unsupported operations. Calibrate the verifier on equal circuits with a global phase change,
swapped wires, an omitted relative phase (a phase change on only part of the state), and a pair
that agrees on all-zero input but differs elsewhere. Convert accepted candidates to the shared
basis without an extra optimization pass. Select the circuit with fewest CX gates, then lowest CX
depth (the number of sequential layers of CX gates, counting one-qubit gates as zero layers).
Break any remaining tie by keeping the earliest available result. The input is available at time
zero. Invalid or late results leave it eligible and count toward the method's failure statistics.

\section{Analysis}

For each of the 48 compiled inputs and each model, record three CX counts.

\begin{itemize}
  \item $Q_i$ is the verified Qiskit level-3 count.
  \item $C_i$ is the best verified count from the extra conventional method, including $Q_i$ as a
  fallback.
  \item $L_{i,m}$ is the best verified count from model $m$, including $Q_i$ as a fallback.
\end{itemize}

The primary result is $C_i-L_{i,m}$ for each circuit and model. Positive values are model wins,
zero is a tie, and negative values are conventional wins. Report wins, ties, and losses, and the
median and total signed CX difference, overall and by circuit group. Also report $Q_i-L_{i,m}$ so
readers can see which model savings are available from another optimizer and which go beyond both
tested tools. Report absolute gate counts, two-qubit depth, one-qubit gate count, verification
failures, timeouts, and time to the best result. Show how many verified gains each method has
delivered by 10, 60, and 180 seconds using the same runs.

For the paired question, subtract the best final CX count the model achieves before level 3 from
its best final CX count after level 3 on the same circuit. Compare final counts, so that both
starts are judged by the length of the resulting circuit. Use the best verified conventional count
from either starting point as a reference for available reductions. A positive difference between
$Q_i$ and that reference means a known reduction exists within the methods tried. This reference
is a practical lower-count implementation and may itself be improved. Compare patterns across
constructions and check whether selected wins recur in the repeat trials.

Publish the per-circuit results and inspect the largest verified reductions alongside ties and
losses of similar starting sizes. Treat closely related widths and modes as variants of a
construction when interpreting patterns. For example, savings on several widths of one
ripple-carry adder would support a finding about that construction. Similar gains on distinct
adder designs would support a broader finding.

\clearpage
\section{Expected Deliverables}

\begin{itemize}
  \item The circuit generator and the gate lists for both starting points.
  \item The exact prompts, compiler settings, and frozen experiment manifest.
  \item The independent verification code and its known-answer checks.
  \item Per-attempt status and timing, reproducible tables and figures, and verified rewrite examples.
\end{itemize}

The accompanying literature review explains the basis for the study and the possible extensions.

\section{Thesis Milestones}

\par\smallskip\noindent
\begin{tabularx}{\textwidth}{@{}
  >{\raggedright\arraybackslash}p{0.32\textwidth}
  >{\raggedright\arraybackslash}X
  @{}}
  \toprule
  \textbf{Weeks} & \textbf{Deliverable} \\
  \midrule
  1--2 &
    Circuit generator, common gate format, reference verifier, and tests for known-answer and
    deliberately incorrect circuits. \\
  \addlinespace[2pt]
  3--4 &
    Conventional adapters, time and failure tracking, prompt checks on the development set, and
    the frozen evaluation manifest. \\
  \addlinespace[2pt]
  5--6 &
    Complete logs for both models and both starting points, including the planned repeat cases. \\
  \addlinespace[2pt]
  7--8 &
    Reproducible per-circuit tables, paired figures, examples of verified rewrites, and a report
    answering where the additional model pass earns its time. \\
  \bottomrule
\end{tabularx}
\par\smallskip

\thesisreferences
\end{document}
